%! AP Statistics — Unit 8, Lesson 8.4 — Group Activity
\documentclass[10pt]{article}
\usepackage{apstats-article}
\usepackage{apstats-boxes}
\newcommand{\TallMath}[1]{$\displaystyle #1\rule[-1.4em]{0pt}{3.2em}$}

\begin{document}

\pageheader{Unit 8, Lesson 8.4}{Group Activity}
\namepartnerperiod

\vspace{0.1in}

\begin{scenariobox}[School Sports Participation Survey]{sky}
A random sample of 150 middle school students is asked about their grade (6th, 7th, 8th)
and whether they participate in a school sport (Yes or No). The results are shown below.

\vspace{0.1in}
\renewcommand{\arraystretch}{1.6}
\begin{tabular}{lcccc}
             & \textbf{Yes} & \textbf{No} & \textbf{Row Total} \\
    \hline
    6th grade & 24 & 26 & 50 \\
    7th grade & 30 & 20 & 50 \\
    8th grade & 18 & 32 & 50 \\
    \hline
    Column Total & 72 & 78 & 150 \\
\end{tabular}
\end{scenariobox}

\begin{tcolorbox}[colback=white, colframe=black!40, title=\textbf{Tier R --- Remediate},
    fonttitle=\bfseries, arc=2mm, left=3mm, right=3mm, top=2mm, bottom=2mm]

\textbf{Task a.} Use the formula $E = \dfrac{(\text{row total})(\text{col total})}{\text{table total}}$
to calculate the expected count for \textbf{two} of the six cells. Show all work.

\begin{enumerate}[label=\alph*.]
    \item 6th grade / Yes: \quad $E = \dfrac{(\blank{1.5cm})(\blank{1.5cm})}{\blank{1.5cm}} = \blank{2cm}$

          \vspace{0.1in}

    \item 7th grade / No: \quad $E = \dfrac{(\blank{1.5cm})(\blank{1.5cm})}{\blank{1.5cm}} = \blank{2cm}$
\end{enumerate}

\vspace{0.1in}
\textbf{Task b.} Use the pattern you see to fill in all 6 expected counts (fill-in provided):

\renewcommand{\arraystretch}{1.8}
\begin{tabular}{lcc}
             & \textbf{Expected: Yes} & \textbf{Expected: No} \\
    \hline
    6th grade & \blank{2.5cm} & \blank{2.5cm} \\
    7th grade & \blank{2.5cm} & \blank{2.5cm} \\
    8th grade & \blank{2.5cm} & \blank{2.5cm} \\
\end{tabular}

\vspace{0.1in}
\textbf{Task c.} Large-Counts Condition: Are all expected counts $\ge 5$?
\quad Circle one: \quad \textbf{Yes} \quad \textbf{No}
\end{tcolorbox}

\begin{tcolorbox}[colback=white, colframe=black!40, title=\textbf{Tier A --- Approaching Proficiency},
    fonttitle=\bfseries, arc=2mm, left=3mm, right=3mm, top=2mm, bottom=2mm]

\textbf{Task a.} Calculate the expected count for \textbf{all 6 cells}. Show work for at least 2 cells.

\vspace{0.1in}

\textbf{Task b.} Complete the side-by-side comparison table:

\renewcommand{\arraystretch}{1.8}
\begin{tabular}{lcccc}
             & \textbf{Obs: Yes} & \textbf{Exp: Yes} & \textbf{Obs: No} & \textbf{Exp: No} \\
    \hline
    6th grade & 24 & \blank{2cm} & 26 & \blank{2cm} \\
    7th grade & 30 & \blank{2cm} & 20 & \blank{2cm} \\
    8th grade & 18 & \blank{2cm} & 32 & \blank{2cm} \\
\end{tabular}

\vspace{0.1in}
\textbf{Task c.} Check the large-counts condition. State it in a complete sentence.

\writelines{2}
\end{tcolorbox}

\begin{tcolorbox}[colback=white, colframe=black!40, title=\textbf{Tier E --- Extension},
    fonttitle=\bfseries, arc=2mm, left=3mm, right=3mm, top=2mm, bottom=2mm]

Complete all tasks from Tier A, then answer:

\textbf{Task d.} Which grade level shows the \textbf{largest difference} between
observed and expected counts? Use numbers to support your answer.

\writelines{2}

\textbf{Task e.} Does the pattern of observed vs.\ expected counts suggest an
\textbf{association} between grade level and sports participation? Explain using
both the numbers from the table and the idea of what expected counts represent.

\writelines{3}

\textbf{Task f.} For the cell with the largest difference, compute
$(O - E)^2/E$. This is one term in the chi-square statistic. What does a
larger value of $(O - E)^2/E$ suggest?

\writelines{3}
\end{tcolorbox}

\end{document}
